\clearpage\thispagestyle{empty}
\pdfbookmark[0]{Quellen zum Bericht über Bessels Arbeit zur Präzession}{besselprecession1814-sources}
\begin{center}\Large\bfseries Quellen zum Bericht über Bessels Arbeit zur Präzession\end{center}

Die Übertragung folgt dem Bericht von Carl Friedrich Gauß in den \emph{Werken}, Bd.~VI (Göttingen, 1874), S.~561. Der Erstdruck erschien in den \emph{Göttingischen gelehrten Anzeigen}, Stück~19 vom 31.~Januar 1814, S.~188--189. Das Digitalisat der Zeitschrift stammt aus der Niedersächsischen Staats- und Universitätsbibliothek Göttingen.

Der Koeffizient in der zweiten Zeile der Tabelle im Sammelband beginnt mit einem c-förmigen Zeichen. Im Erstdruck beginnt er mit einer Null, auf die ein Dezimalkomma folgt.

Im Vorwort zu seiner \emph{Untersuchung der Größe und des Einflusses des Vorrückens der Nachtgleichen} (Berlin, Realschulbuchhandlung, 1815) erläutert Friedrich Wilhelm Bessel, dass er seine Preisschrift nach Erscheinen einer neuen Ausgabe von Piazzis Katalog überarbeitete. Auf S.~25 unterscheidet er die 4585 Vergleiche mit dem älteren Katalog von den 4707 mit dem neueren. Auf S.~26 stehen die überarbeiteten Ergebnisse.

KI-Übertragung des deutschen Originals in heutiges Deutsch.

Erstdruck (\href{https://images.p.sub.uni-goettingen.de/iiif/image/gdz:PPN319721507_1814:00000190/full/1650,/0/default.jpg}{188}, \href{https://images.p.sub.uni-goettingen.de/iiif/image/gdz:PPN319721507_1814:00000191/full/1650,/0/default.jpg}{189}).

Bessels Preisschrift (\href{https://www.digitale-sammlungen.de/de/view/bsb10049066?page=3}{Bayerische Staatsbibliothek, 4 Astr.p. 43 m}, \href{https://api.digitale-sammlungen.de/iiif/image/v2/bsb10049066_00003/full/1700,/0/default.jpg}{Titelblatt}, \href{https://api.digitale-sammlungen.de/iiif/image/v2/bsb10049066_00005/full/1700,/0/default.jpg}{Vorwort}, \href{https://api.digitale-sammlungen.de/iiif/image/v2/bsb10049066_00027/full/1700,/0/default.jpg}{25}, \href{https://api.digitale-sammlungen.de/iiif/image/v2/bsb10049066_00028/full/1700,/0/default.jpg}{26}).

\clearpage
